\subsection{Maximal Commutative Subalgebras}

We say that a subalgebra of a K-algebra A is a maximal commutative subalgebra of A if it is a maximal element of the set of commutative subalgebras of A.

Lemma 3. --- Let A be a K-algebra and L a subalgebra of A.

\begin{enumerate}
\def\labelenumi{\alph{enumi})}
\item
  The algebra \(\mathrm{L}\) is a maximal commutative subalgebra of \(\mathrm{A}\) if and only if \(\mathrm{L}\) is equal to its commutant \(\mathrm{L}'\) in \(\mathrm{A}\) .
\item
  Let \(\mathrm{K}'\) be a nonzero commutative K-algebra. Then \(\mathrm{L}\) is a maximal commutative subalgebra of A if and only if \(\mathrm{L}_{(\mathrm{K}')}\) is a maximal commutative subalgebra of \(\mathrm{A}_{(\mathrm{K}')}\) .
\end{enumerate}

Let us prove a). First suppose that L is equal to \(L'\) . Then L is commutative; if M is a commutative subalgebra of A containing L, then we have \(xy = yx\) for \(x\) in \(\mathrm{L}\) and \(y\) in \(\mathrm{M}\) , so \(\mathrm{M} \subset \mathrm{L}'\) and therefore \(\mathrm{M} = \mathrm{L}\) . Consequently, \(\mathrm{L}\) is a maximal commutative subalgebra of \(\mathrm{A}\) .

Conversely, suppose that L is a maximal commutative subalgebra of A, and let x be an element of \(L'\) . The subalgebra M of A generated by \(L \cup \{x\}\) is then commutative and contains L. Because L is maximal, we have M = L, and so \(x \in L\) and finally \(L = L'\) , which gives a).

By Proposition 6 of III, §4, No.~4, p.~468, the commutant of \(\mathrm{L}_{(\mathrm{K}^{\prime})}\) in \(\mathrm{A}_{(\mathrm{K}^{\prime})}\) is \(\mathrm{L}'_{(\mathrm{K}^{\prime})}\) . Since the equalities \(\mathrm{L} = \mathrm{L}'\) and \(\mathrm{L}_{(\mathrm{K}^{\prime})} = \mathrm{L}'_{(\mathrm{K}^{\prime})}\) are equivalent (II, §7, No.~9, p.~311, Proposition 19), assertion b) follows from a).

PROPOSITION 3. --- Let A be a central simple K-algebra of finite degree, and let L be a semisimple commutative subalgebra of A. The following properties are equivalent:

\begin{enumerate}
\def\labelenumi{(\roman{enumi})}
\item
  The algebra \(\mathrm{L}\) is a maximal commutative subalgebra of \(\mathrm{A}\) .
\item
  The left L-module A is free, of dimension {[}L : K{]}.
\item
  We have \([A:K]=[L:K]^{2}\) .
\end{enumerate}

Suppose, moreover, that A is the algebra \(\operatorname{End}_{\mathrm{K}}(\mathrm{V})\) , where V is a vector space of nonzero finite dimension over K. Then the previous properties are also equivalent to the following:

\begin{enumerate}
\def\labelenumi{(\roman{enumi})}
\setcounter{enumi}{3}
\tightlist
\item
  V is a free L-module of dimension 1.
\end{enumerate}

\begin{enumerate}
\def\labelenumi{\Alph{enumi})}
\tightlist
\item
  Suppose that A is of the form \(\operatorname{End}_{\mathrm{K}}(\mathrm{V})\) , where V is a vector space of nonzero finite dimension over the field K. We will establish the equivalence of conditions (i) through (iv) following the diagram
\end{enumerate}

\[
(i) \Longrightarrow (i v) \Longrightarrow (i i) \Longrightarrow (i i i) \Longrightarrow (i).
\]

Since L is a semisimple commutative algebra of finite degree over K, we can identify it with a finite product \(\prod_{i\in I}L_{i}\) of extensions of K of finite degree (VIII, p.~137, Proposition 3). For every \(i\in I\) , let \(V_{i}\) be the isotypical component of type \(L_{i}\) of the L-module V; it is a vector space of nonzero finite dimension over \(L_{i}\) because V is a faithful L-module (VIII, p.~144, Corollary). We can then identify V with \(\prod_{i\in I}V_{i}\) . Under these conditions, the commutant \(L'\) of L in A, which is simply the algebra \(\operatorname{End}_{\mathrm{L}}(\mathrm{V})\) , can be identified with the product \(\prod_{i\in I}\operatorname{End}_{\mathrm{L}_{i}}(\mathrm{V}_{i})\) .

Suppose that L is a maximal commutative subalgebra of \(\operatorname{End}_{\mathrm{K}}(\mathrm{V})\) . By Lemma 3, a), we have \(L = L'\) , so \(L'\) is commutative and we have \(\dim_{L_i}(V_i) = 1\) for every \(i \in I\) . So (i) implies (iv).

Suppose that the L-module V is free of dimension 1. Let \((e_{1},\ldots,e_{r})\) be a basis of V over K. The mapping \(a\mapsto(ae_{1},\ldots,ae_{r})\) is an isomorphism of left A-modules and therefore of L-modules from A to \(V^{r}\) . Consequently, A is a free left L-module of dimension \(r\) , and we have \(r = \dim_{\mathrm{K}}(\mathrm{V}) = [\mathrm{L}:\mathrm{K}]\) . So (iv) implies (ii).

It is clear that (i) implies (iii).

Finally, suppose that we have \([A:K]=[L:K]^{2}\) , in other words, \(\dim_{\mathrm{K}}(\mathrm{V})=[\mathrm{L}:K]\) . We then have \(\sum_{i}\dim_{\mathrm{L}_{i}}(\mathrm{V}_{i})[\mathrm{L}_{i}:\mathrm{K}]=\sum_{i}[\mathrm{L}_{i}:\mathrm{K}]\) , so that \(V_{i}\) has dimension 1 over \(L_{i}\) for every i. We then have \(\operatorname{End}_{\mathrm{L}_{i}}(\mathrm{V}_{i})=\mathrm{L}_{i}\) for every i, and so \(L^{\prime}=L\) . By Lemma 3, a), L is a maximal commutative subalgebra of A. So (iii) implies (i).

\begin{enumerate}
\def\labelenumi{\Alph{enumi})}
\setcounter{enumi}{1}
\tightlist
\item
  Let us continue with the general case. By Theorem 1 of VIII, p.~252, there exists a separable extension \(\mathrm{K}'\) of \(\mathrm{K}\) of finite degree such that the \(\mathrm{K}'\) -algebra \(\mathrm{A}_{(\mathrm{K}')}\) is isomorphic to an algebra \(\mathrm{End}_{\mathrm{K}'}(\mathrm{V}')\) , where \(\mathrm{V}'\) is a vector space of nonzero finite dimension over \(\mathrm{K}'\) . Then the \(\mathrm{K}'\) -algebra \(\mathrm{L}_{(\mathrm{K}')}\) is commutative and semisimple (VIII, p.~222, Corollary 2). By the first part of the proof, the following properties are equivalent:
\end{enumerate}

(i') The algebra \(L_{(K')}\) is a maximal commutative subalgebra of \(A_{(K')}\) .

(ii') The left \(\mathrm{L}_{(\mathrm{K}^{\prime})}\) -module \(\mathrm{A}_{(\mathrm{K}^{\prime})}\) is free, of dimension \([\mathrm{L}_{(\mathrm{K}^{\prime})}:\mathrm{K}^{\prime}]\) .

(iii') We have \(\left[\mathrm{A}_{\left(\mathrm{K}^{\prime}\right)}:\mathrm{K}^{\prime}\right] = \left[\mathrm{L}_{\left(\mathrm{K}^{\prime}\right)}:\mathrm{K}^{\prime}\right]^{2}\) .

By Lemma 3, b), properties (i) and (i') are equivalent.

Set \(n = [\mathrm{L}:\mathrm{K}]\) ; then \(n = [\mathrm{L}_{(\mathrm{K}^{\prime})}:\mathrm{K}^{\prime}]\) . Property (ii) means that the left L-modules A and \(\mathrm{L}^n\) are isomorphic; by Theorem 3 of VIII, p.~37, this is equivalent to the \(\mathrm{L}_{(\mathrm{K}^{\prime})}\) -modules \(\mathrm{A}_{(\mathrm{K}^{\prime})}\) and \((\mathrm{L}_{(\mathrm{K}^{\prime})})^n\) being isomorphic. The equivalence of (ii) and (ii') follows.

Finally, we have \([A:K]=[A_{(K')}:K']\) and \([L:K]=[L_{(K')}:K']\) , which give the equivalence of properties (iii) and (iii'). We have thus proved the equivalence of properties (i), (ii), and (iii).

COROLLARY. --- Let A be a central simple algebra of finite degree over K, and let L be a semisimple commutative K-algebra such that \([A:K]\) is equal to \([L:K]^{2}\) . Let f and g be injective homomorphisms from L to A. There exists an inner automorphism \(\theta\) of A such that \(g = \theta \circ f\) .

Set \(n = [L : K]\) . Viewed as a left module over the subring \(f(L)\) , A is free of dimension n; this follows from the equivalence of properties (ii) and (iii) of Proposition 3. Since f is an isomorphism from L to \(f(L)\) , the left L-module \(A^{f}\) (whose law of action is given by \((x, a) \mapsto f(x)a\) ) is free of dimension n.~The same holds for \(A^{g}\) , which is therefore isomorphic to \(A^{f}\) . We conclude using the equivalence of properties (i) and (ii) of Proposition 1 (VIII, p.~257).

Suppose that \(\mathbf{A}\) is a central simple algebra of finite degree over \(\mathbf{K}\) .

There can exist maximal commutative subalgebras L of A that are not semisimple and for which \([A:K]\neq[L:K]^{2}\) (VIII, p.~270, Exercise 5).

